\documentclass[11pt]{article}
\usepackage{estudio}
\cuadernillo{3 --- Simulaciones}

\begin{document}

\portada{3}{Simulaciones}{%
  Diapositivas 8 a 10 \textbullet\ expositor C \textbullet\ 1:35\\[4pt]
  Sistema simulado, observables y criterio de medición}

\tableofcontents
\clearpage

% ===========================================================================
\section{Qué se cuenta acá}

La guía de la cátedra (punto 2.3) pide cuatro cosas en esta sección:

\begin{enumerate}
  \item Descripción del sistema \emph{particular}: geometría, parámetros fijos,
    inputs y outputs. Con un esquema.
  \item Definición \textbf{matemática} de los observables. Si algo se promedia,
    hay que explicitar qué se suma y sobre qué se divide.
  \item Número de repeticiones.
  \item Tiempos de las simulaciones.
\end{enumerate}

Las diapositivas 9 y 10 están armadas exactamente sobre esos cuatro puntos.

\clave{Esta es la sección donde la cátedra verifica que sepas \emph{qué medís y
cómo}. Es la que tiene más preguntas posibles por minuto de exposición, porque
cada número de la diapositiva 10 es una decisión metodológica. Estudiá bien la
sección de criterio de estacionario: es lo que más se pregunta.}

% ===========================================================================
\section{Guion: diapositivas 8 a 10}

\diapo{8}{Divisoria: Simulaciones}{2 s}{C}

Tu señal de entrada.

\diapo{9}{Sistema simulado}{38 s}{C}

A la izquierda el esquema de la caja con las partículas, sus vectores, el radio
$r_c$ y las flechas de contorno periódico. A la derecha los parámetros.

\begin{quote}
\itshape
``El sistema es una caja cuadrada de lado 10 con condiciones periódicas de
contorno, o sea que lo que sale por un lado entra por el opuesto: no hay
paredes. El radio de interacción es 1 y el paso temporal es 1. La rapidez es
$3 \times 10^{-2}$, elegida de forma que en un paso una partícula recorra el
3\% de su radio de interacción, así el vecindario cambia de a poco. El input que
barremos es el ruido, entre 0 y $2\pi$, que es todo el rango posible. Y el
número de partículas sale de la densidad, como $\rho$ por $L$ al cuadrado:
estudiamos las tres densidades del enunciado, 2, 4 y 8, y agregamos tres
densidades mucho más bajas para el estudio de clusters, de modo que $N$ va de 11
a 800.''
\end{quote}

\ojo{No leas la lista de parámetros número por número: están escritos, la
cátedra los lee. Lo que tenés que aportar con la voz es \emph{por qué} son esos:
que $v\,\Delta t$ es el 3\% de $r_c$, que $[0, 2\pi]$ cubre todo el ruido
posible, y que hay tres densidades extra que no están en el enunciado.}

\diapo{10}{Observables}{55 s}{C}

Es la diapositiva más importante de tu primer bloque y la más densa: dos
definiciones matemáticas más el bloque de metodología.

\begin{quote}
\itshape
``Medimos dos observables. El primero es la polarización: el módulo de la suma
vectorial de todas las velocidades, dividido por $N$ y por la rapidez. Está
normalizado entre 0 y 1: vale 1 si todas las partículas apuntan exactamente en
la misma dirección, y tiende a 0 si las direcciones se reparten de forma
isótropa. Es el parámetro de orden. El segundo es la fracción de la componente
gigante: armamos el grafo donde cada partícula es un nodo y hay una arista entre
las que están a distancia menor que $r_c$ ---exactamente la misma relación de
vecindad que usa la dinámica---, buscamos la componente conexa más grande y
reportamos qué fracción del total de partículas contiene. Los dos observables
son series temporales, y el escalar que reportamos es su promedio en el estado
estacionario. Cada punto se corre con 10 realizaciones independientes; las
corridas son de $2 \times 10^3$ pasos para las densidades del enunciado y de
$1 \times 10^4$ para las subcríticas, que tardan más en estacionarse. La ventana
de promedio no es fija: se mide para cada nivel de ruido.''
\end{quote}

\clave{Las tres cosas que no podés omitir: \textbf{(1)} $v_a$ está normalizado
entre 0 y 1 y es el parámetro de orden; \textbf{(2)} el grafo de clusters usa
\emph{la misma} relación de vecindad que la dinámica, no otra; \textbf{(3)} la
ventana de promedio se \emph{mide}, no se fija de antemano.}

Cerrá pasándole la palabra a A:

\begin{quote}
\itshape
``Con eso definido, [A] arranca con los resultados del modelo de Vicsek.''
\end{quote}

% ===========================================================================
\section{El fondo: cada número de las diapositivas 9 y 10}

\subsection{Los parámetros geométricos}

\begin{center}
\begin{tabular}{llp{7.6cm}}
\toprule
Símbolo & Valor & Por qué \\
\midrule
$L$        & 10                & Lo fija el enunciado. \\
$r_c$      & 1                 & Radio de interacción; fija la escala local. \\
$\Delta t$ & 1                 & Unidad de tiempo; el paso del autómata. \\
$v$        & $3\times10^{-2}$  & $v\,\Delta t = 0{,}03 = 3\%$ de $r_c$. \\
Contorno   & periódico         & Evita bordes: no hay partículas privilegiadas. \\
$M$        & 9                 & Máximo con lado de celda $>r_c$ (cuadernillo 2). \\
\bottomrule
\end{tabular}
\end{center}

\teoria{La elección de $v$ es la que más se puede preguntar. Si $v\,\Delta t$
fuera comparable a $r_c$, una partícula cambiaría por completo de vecindario en
un solo paso y la ``interacción local'' dejaría de ser local: el sistema se
volvería un campo medio. Si fuera muchísimo menor, las partículas casi no se
mezclarían y el modelo tendería al XY en red, que en 2D no ordena. El valor
$0{,}03\,r_c$ está en el régimen intermedio que usa la bibliografía: hay
mezcla, pero el vecindario cambia gradualmente.}

\subsection{Las seis densidades y por qué hay tres de más}

$N = \rho L^2$, y el número medio de vecinos geométricos es
$\langle k \rangle = \rho\,\pi r_c^2$.

\begin{center}
\begin{tabular}{lccl}
\toprule
$\rho$ & $N$ & $\langle k \rangle$ & Para qué \\
\midrule
8                   & 800 & 25{,}1 & Enunciado, todos los incisos \\
4                   & 400 & 12{,}6 & Enunciado, todos los incisos \\
2                   & 200 & \phantom{0}6{,}3 & Enunciado, todos los incisos \\
\midrule
\multicolumn{2}{l}{\emph{umbral de percolación}} & \phantom{0}4{,}5 & \\
\midrule
$1/\pi \approx 0{,}318$    & 32 & 1{,}00 & Estudio de clusters \\
$1/(2\pi) \approx 0{,}159$ & 16 & 0{,}50 & Estudio de clusters \\
$1/(3\pi) \approx 0{,}106$ & 11 & 0{,}33 & Estudio de clusters \\
\bottomrule
\end{tabular}
\end{center}

\dato{Las tres densidades del enunciado tienen $\langle k \rangle$ = 6,3, 12,6 y
25,1: \textbf{todas por encima del umbral de percolación} $\approx 4{,}5$. Eso
significa que la red de vecinos percola \emph{sin importar} el grado de orden
direccional, y por eso $S \approx 1$ en todo el rango de $\eta$. Las tres
densidades subcríticas se agregaron ---por indicación de la cátedra--- para que
el observable $S$ tenga algo que mostrar. Es el resultado que explica A en la
diapositiva 16 y vos en la 23.}

\subsection{La grilla de ruido}

39 valores de $\eta$ en $[0,\ 2\pi]$, \textbf{la misma para los dos modelos}:

\begin{itemize}
  \item base uniforme de paso $0{,}2$, más el extremo $2\pi$;
  \item refinamiento de paso $0{,}05$ en $[0,\ 0{,}4]$.
\end{itemize}

El refinamiento en ruido bajo no es decorativo: el votante colapsa ahí. Su cruce
de $v_a = 0{,}5$ cae entre $0{,}19$ y $0{,}37$ según la densidad, o sea dentro
de los dos primeros puntos de la grilla base. Sin refinar, la transición del
votante quedaría sin resolver. Vicsek, en cambio, transiciona cerca de
$\eta \approx 2{,}6$--$3{,}4$, donde el paso $0{,}2$ alcanza de sobra.

\clave{Que la grilla sea \emph{idéntica} para los dos modelos es lo que hace que
la comparación del inciso (f) sea punto a punto. Si cada modelo tuviera su
propia grilla, las curvas superpuestas no compartirían abscisas y la
comparación no sería legítima.}

En total: 39 valores de $\eta$ $\times$ 6 densidades $\times$ 2 modelos =
\textbf{468 puntos de barrido}, cada uno con 10 semillas: 4680 corridas.

\subsection{Las 10 realizaciones y de dónde salen las semillas}

Cada punto (modelo, $\rho$, $\eta$) se corre con $K = 10$ semillas
independientes. La semilla \textbf{nunca} se toma del reloj: se deriva como los
64 bits más significativos de
$\mathrm{sha256}(\text{modelo}\,|\,\rho\,|\,\eta\,|\,k)$, con los números
formateados a precisión fija.

Dos ventajas:

\begin{itemize}
  \item \textbf{Reproducibilidad}: todo el barrido se puede regenerar bit a bit.
  \item \textbf{Descorrelación}: dos valores de $\eta$ arbitrariamente próximos
    producen semillas sin ninguna relación entre sí. Si las semillas fueran, por
    ejemplo, correlativas, puntos vecinos de la curva compartirían fluctuaciones
    y las barras de error subestimarían la incerteza real.
\end{itemize}

\subsection{Los largos de corrida: 2000 y 10\,000 pasos}

No es una constante única, y hay que saber por qué. El tiempo que tarda el
sistema en llegar al estacionario \emph{crece al bajar $N$}. Midiendo la
duración del transitorio con 20 semillas:

\begin{itemize}
  \item a $\rho = 2$ ($N = 200$) el observable ya es estacionario en $t = 1000$;
  \item a $\rho = 1/\pi$ y $1/(3\pi)$ ($N = 32$ y $11$) sigue evolucionando
    hasta $t \approx 4000$.
\end{itemize}

Por eso las densidades del enunciado se corren $2\times10^3$ pasos y las
subcríticas $1\times10^4$: con 10\,000 pasos, la ventana de promedio cae
enteramente dentro del estacionario.

\subsection{El criterio de estado estacionario (la pregunta más probable)}

Es la parte metodológicamente más fina del trabajo. El escalar que reportamos es
el promedio de la serie temporal sobre una ventana, y esa ventana \textbf{se
mide para cada punto de barrido}, no se fija de antemano.

\textbf{Por qué no fijarla.} La duración del transitorio depende fuertemente del
ruido. A $\eta$ grande el sistema nace ya desordenado y el estacionario arranca
en $t = 0$: no hay nada que descartar. A $\eta$ chico la polarización tarda
$\sim 10^2$ pasos en llegar a su meseta. Un corte fijo tendría que valer el peor
caso de todos, y estaría tirando datos útiles en la gran mayoría de los puntos.

\textbf{Cómo se mide.} El procedimiento, idéntico para $v_a$ y para $S$:

\begin{enumerate}
  \item Se promedian punto a punto las $K = 10$ series del punto de barrido.
  \item Esa curva de ensamble se parte en \textbf{40 bloques} y se promedia cada
    uno.
  \item La meseta se estima con los bloques de la \emph{segunda mitad},
    estacionaria por construcción: da un valor medio $\mu$ y una dispersión
    entre bloques $\sigma$.
  \item Se define una banda de tolerancia
    $\max(0{,}05\,|\mu|,\ 3\sigma)$ alrededor de $\mu$.
  \item El corte es el final del \textbf{último} bloque que se sale de la banda:
    o sea, el primer instante a partir del cual la curva ya no vuelve a
    apartarse.
  \item Tope duro: el corte nunca cae después de la \textbf{mitad} de la
    corrida.
\end{enumerate}

\teoria{Los dos detalles que hacen que esto funcione. \textbf{Primero}, por qué
se promedia entre semillas antes de detectar: una corrida individual tiene
fluctuaciones lentas que \emph{no} son transitorio ---la serie se aparta y
vuelve al mismo valor---; el promedio de ensamble las cancela y deja solo la
relajación sistemática desde la condición inicial. \textbf{Segundo}, por qué la
banda es el máximo de dos términos: el término relativo ($5\%$) manda cuando la
serie es limpia, y el de $3\sigma$ manda cerca de la transición, donde la
fluctuación es grande y un criterio relativo cortaría siempre al tope.}

\textbf{Ventanas separadas para $v_a$ y $S$.} Cada observable recibe su propia
ventana, con el mismo criterio. No comparten corte porque no comparten
transitorio: $S$ ya arranca cerca de su valor estacionario en la condición
inicial ---las partículas nacen conectadas--- mientras que $v_a$ tiene que
relajar desde ángulos al azar. Forzar una ventana única significaría promediar
uno de los dos sobre menos datos de los que tiene.

\clave{Los cortes se guardan en un archivo y son \emph{los mismos} que dibujan
las líneas verticales de las figuras de evolución temporal (diapositivas 13, 15,
18, 20). O sea: cada figura muestra literalmente la ventana sobre la que se
promedió. Eso es exactamente lo que pide el enunciado en el inciso (b), ``mostrar
con líneas verticales el inicio del mismo''.}

\subsection{Qué son las barras de error}

Son el \textbf{desvío estándar de las 10 medias de estado estacionario}, una por
semilla. Es decir: cuantifican la incerteza \emph{de la media entre
realizaciones}, no la amplitud de la fluctuación temporal del observable dentro
de una corrida.

\ojo{Las dos magnitudes son muy distintas. En Vicsek a $\rho = 2$, cerca de la
transición, la fluctuación temporal dentro de una corrida es del orden de
$0{,}12$, mientras que el desvío entre semillas es $\approx 0{,}014$: casi diez
veces menor. Si te preguntan ``¿qué representa la barra de error?'', la
respuesta es \emph{la dispersión entre realizaciones independientes}, y conviene
que lo digas con esas palabras.}

\subsection{La condición inicial}

Posiciones $x$ e $y$ sorteadas uniformemente en $[0, L)$ de forma independiente,
y orientación uniforme en $[-\pi, \pi]$. No hay criterio de no solapamiento
---a diferencia del generador del TP1--- porque acá las partículas son
puntuales.

El sistema parte entonces espacialmente homogéneo y direccionalmente isótropo,
lo que implica $v_a(0) \approx N^{-1/2}$.

\subsection{El piso de tamaño finito: $N^{-1/2}$}

Este número aparece por todos lados en los resultados y conviene entenderlo
ahora. En la fase desordenada las $N$ direcciones son esencialmente
independientes, así que la suma vectorial es una caminata aleatoria de $N$ pasos
unitarios: su módulo típico es $\sqrt{N}$, y al dividir por $N$ queda
$v_a \approx N^{-1/2}$.

\begin{center}
\begin{tabular}{lcccccc}
\toprule
$N$ & 11 & 16 & 32 & 200 & 400 & 800 \\
\midrule
$N^{-1/2}$ & 0{,}30 & 0{,}25 & 0{,}18 & 0{,}071 & 0{,}050 & 0{,}035 \\
$v_a$ medido en $\eta = 2\pi$ & 0{,}27 & 0{,}22 & 0{,}16 & 0{,}063 & 0{,}044 & 0{,}031 \\
\bottomrule
\end{tabular}
\end{center}

\dato{El acuerdo es excelente en las seis densidades. O sea: \textbf{$v_a$ nunca
llega a cero, y no tiene por qué llegar}. El valor de saturación a ruido máximo
no es un resultado del modelo, es un efecto de tamaño finito perfectamente
entendido. Si alguien pregunta ``¿por qué la curva no baja a cero?'', esta tabla
es la respuesta.}

% ===========================================================================
\section{Preguntas y respuestas}

\subsection{Sobre los parámetros}

\pregunta{¿Por qué $v = 0{,}03$?}
\respuesta{Para que en un paso una partícula recorra el 3\% de su radio de
interacción, o sea que el vecindario cambie gradualmente. Si $v\,\Delta t$ fuera
comparable a $r_c$, la partícula cambiaría de vecindario entero en un paso y la
interacción dejaría de ser local. Si fuera muchísimo menor, no habría mezcla y
el sistema tendería al modelo XY en red, que en dos dimensiones no ordena.}

\pregunta{¿Por qué $\Delta t = 1$?}
\respuesta{Es la unidad de tiempo del autómata: un paso. Todas las magnitudes
son adimensionales, así que fijar $\Delta t = 1$ solo define la escala temporal;
lo que tiene significado físico es el producto $v\,\Delta t$ relativo a $r_c$.}

\pregunta{¿Por qué condiciones periódicas y no paredes?}
\respuesta{Para que no haya partículas privilegiadas. Con paredes, las del borde
tendrían menos vecinos y un comportamiento distinto, y aparecería un efecto de
superficie que contaminaría la medición del orden global. Las condiciones
periódicas hacen que todas las partículas sean equivalentes.}

\pregunta{¿De dónde salen las densidades $1/\pi$, $1/2\pi$ y $1/3\pi$? No están
en el enunciado.}
\respuesta{Las agregamos por indicación de la cátedra, para el estudio de
clusters. Están elegidas para que $\langle k \rangle = \rho\pi r_c^2$ dé
exactamente 1, 0,5 y 0,33, o sea holgadamente por debajo del umbral de
percolación $\approx 4{,}5$. En las tres densidades del enunciado
$\langle k \rangle$ vale 6,3, 12,6 y 25,1, todas por encima del umbral, y por
eso ahí $S \approx 1$ para todo $\eta$ y el observable no discrimina.}

\pregunta{¿Por qué barren $\eta$ hasta $2\pi$?}
\respuesta{Porque con $\eta = 2\pi$ el ruido ya es uniforme sobre toda la
circunferencia: la dirección nueva es completamente aleatoria. Es el desorden
máximo alcanzable y más allá no cambia nada.}

\pregunta{¿Por qué la grilla de $\eta$ no es uniforme?}
\respuesta{La base es uniforme de paso 0,2, y le agregamos un refinamiento de
paso 0,05 entre 0 y 0,4 porque ahí colapsa el votante: su transición cae entre
0,19 y 0,37 según la densidad, o sea dentro de los dos primeros puntos de la
grilla base. Sin refinar, la transición del votante quedaría sin resolver.
Vicsek transiciona cerca de 2,6--3,4, donde el paso 0,2 alcanza.}

\pregunta{¿Usan la misma grilla de $\eta$ para los dos modelos?}
\respuesta{Sí, exactamente la misma, 39 valores. Es lo que hace que la
comparación entre modelos sea punto a punto: las curvas superpuestas comparten
abscisas.}

\subsection{Sobre los observables}

\pregunta{¿Por qué $v_a$ está dividido por $N$ y por $v$?}
\respuesta{Para normalizarlo entre 0 y 1. El módulo de la suma de las $N$
velocidades vale como máximo $N v$, que se alcanza cuando todas apuntan igual;
dividiendo por $Nv$ el observable queda adimensional y comparable entre
densidades distintas.}

\pregunta{¿Por qué $v_a$ no llega a cero a ruido máximo?}
\respuesta{Por tamaño finito. En la fase desordenada las direcciones son
independientes, así que la suma vectorial es una caminata aleatoria de $N$ pasos
unitarios y su módulo típico es $\sqrt{N}$: al dividir por $N$ queda
$v_a \approx N^{-1/2}$. Para $N = 200$ eso da 0,071 y medimos 0,063; para
$N = 800$ da 0,035 y medimos 0,031. El acuerdo es muy bueno en las seis
densidades.}

\pregunta{¿Cómo definen un cluster?}
\respuesta{Como una componente conexa del grafo donde los nodos son las
partículas y hay arista entre las que están a distancia menor que $r_c$. Es
\emph{exactamente} la misma relación de vecindad que usa la dinámica: el
observable se calcula recorriendo la misma lista de adyacencia que produjo el
Cell Index Method, no se vuelve a buscar vecinos ni se hace fuerza bruta
aparte. Así la definición de cluster coincide por construcción con la de
vecindad.}

\pregunta{¿Cómo calculan la componente conexa más grande?}
\respuesta{Recorrido en profundidad con pila explícita y vector de visitados,
sobre la lista de adyacencia. Se recorre cada componente una vez y se queda el
máximo tamaño, dividido por $N$.}

\pregunta{Con contorno periódico, ¿un cluster puede darle la vuelta a la caja?}
\respuesta{Sí, y no le damos ningún tratamiento especial: forma parte del
fenómeno. Lo único que pasa es que al visualizarlo, el cluster aparece partido
en los bordes de la caja, pero en el grafo es una sola componente.}

\subsection{Sobre la metodología de medición}

\pregunta{¿Cómo determinan cuándo empieza el estado estacionario?}
\respuesta{Se mide para cada punto de barrido, no se fija. Promediamos las 10
series del punto entre sí, partimos esa curva de ensamble en 40 bloques,
estimamos la meseta con los bloques de la segunda mitad, y definimos una banda
de tolerancia igual al máximo entre el 5\% del valor de la meseta y 3 desvíos
entre bloques. El corte es el final del último bloque que se sale de esa banda,
o sea el primer instante a partir del cual la curva ya no vuelve a apartarse.
Además hay un tope: nunca cae después de la mitad de la corrida.}

\pregunta{¿Por qué promedian entre semillas antes de detectar el corte?}
\respuesta{Porque una corrida individual tiene fluctuaciones lentas que no son
transitorio: la serie se aparta y vuelve al mismo valor. Si detectáramos sobre
una sola corrida, esas excursiones se leerían como transitorio y el corte
saldría tardísimo. El promedio de ensamble las cancela y deja solo la relajación
sistemática desde la condición inicial.}

\pregunta{¿Por qué la banda es el máximo entre dos criterios?}
\respuesta{Porque los dos regímenes necesitan cosas distintas. Cuando la serie
es limpia manda el criterio relativo del 5\%. Cerca de la transición, donde la
fluctuación es grande, un criterio relativo cortaría siempre en el tope, y ahí
manda el de 3 desvíos entre bloques.}

\pregunta{¿Por qué $v_a$ y $S$ tienen ventanas distintas?}
\respuesta{Porque no comparten transitorio. $S$ ya arranca cerca de su valor
estacionario en la condición inicial, porque las partículas nacen conectadas;
$v_a$ tiene que relajar desde una configuración de ángulos al azar. Forzar una
ventana única significaría promediar uno de los dos sobre menos datos de los que
tiene disponibles. El criterio aplicado es idéntico, lo que cambia es el
resultado.}

\pregunta{¿Qué representan las barras de error?}
\respuesta{El desvío estándar de las 10 medias de estado estacionario, una por
semilla. O sea la dispersión entre realizaciones independientes, no la
fluctuación temporal dentro de una corrida. Las dos magnitudes difieren mucho:
cerca de la transición, en Vicsek a $\rho = 2$, la fluctuación temporal es
$\approx 0{,}12$ y el desvío entre semillas $\approx 0{,}014$.}

\pregunta{¿Por qué 10 realizaciones y no más?}
\respuesta{Es un compromiso con el volumen del barrido: 468 puntos por 10
semillas son 4680 corridas. Con 10 semillas las barras de error son chicas fuera
de la región de transición; donde la dispersión crece es justamente cerca del
punto crítico, que es lo esperable porque ahí es donde el sistema fluctúa más.}

\pregunta{¿Por qué las corridas subcríticas son de 10\,000 pasos y las otras de
2000?}
\respuesta{Porque el tiempo hasta el estacionario crece al bajar $N$. Lo
medimos con 20 semillas: a $\rho = 2$ el observable ya es estacionario en
$t = 1000$, mientras que a $\rho = 1/\pi$ y $1/(3\pi)$ sigue evolucionando hasta
$t \approx 4000$. Con 10\,000 pasos la ventana de promedio cae enteramente en el
estacionario.}

\pregunta{¿Cómo eligen las semillas? ¿Son aleatorias?}
\respuesta{No, son deterministas: cada semilla es el hash SHA-256 de la clave
formada por modelo, densidad, ruido e índice de repetición, truncado a 64 bits.
Nunca se siembra desde el reloj. Eso da dos cosas: el barrido completo es
reproducible bit a bit, y dos valores de $\eta$ arbitrariamente próximos
producen semillas sin relación entre sí, así que no hay correlaciones espurias
entre puntos vecinos de la curva.}

\pregunta{¿Cuál es la condición inicial?}
\respuesta{Posiciones uniformes e independientes en la caja y orientaciones
uniformes en $[-\pi,\pi]$. No hay criterio de no solapamiento porque las
partículas son puntuales. El sistema parte espacialmente homogéneo y
direccionalmente isótropo, con $v_a(0) \approx N^{-1/2}$.}

\subsection{Sobre el output (fuera del guion, pero lo pueden preguntar)}

\pregunta{¿En qué formato sale la simulación?}
\respuesta{Texto plano, en dos archivos independientes. Uno es la trayectoria
completa: por cada paso, una línea con el tiempo y después $N$ líneas con
posición y velocidad. El otro es un registro escalar con una línea por paso, con
el tiempo, $v_a$ y $S$; es el que consume el barrido, porque volcar la
trayectoria completa de miles de corridas sería inmanejable.}

\pregunta{¿La animación se genera durante la simulación?}
\respuesta{No, y es un requisito explícito del enunciado. La simulación escribe
texto plano y el módulo de animación se ejecuta después, como un proceso
independiente que toma ese texto como entrada. Así la velocidad de la animación
no queda supeditada a la velocidad de la simulación.}

\end{document}
